\documentclass{article}
\usepackage{import}
\usepackage{tcolorbox}
\usepackage{amsmath}
\usepackage{amssymb}
\usepackage{amsthm}
\usepackage{geometry}
\usepackage{enumitem}
\usepackage{pdfpages}
\usepackage{transparent}
\usepackage{booktabs}
\usepackage{bm}
\usepackage{tabularx}
\usepackage{multirow}
% \usepackage{minted}
\usepackage{xcolor}
\usepackage{setspace}
\usepackage[hidelinks]{hyperref}
\usepackage{pgfplots}
% \usepackage{soulpos}

\tcbuselibrary{theorems}
\tcbuselibrary{skins}

\setcounter{tocdepth}{6}
\setcounter{secnumdepth}{6}

\newcommand{\param}{\boldsymbol{\theta}}
\newcommand{\bs}{\boldsymbol}
\newcommand{\data}{\mathcal D}

% pgfplots settings
\pgfplotsset{compat=1.18}
\usepgfplotslibrary{external}

% tabularx Settings
\renewcommand\tabularxcolumn[1]{>{\centering\arraybackslash}m{#1}}

% Page Settings
\newgeometry{vmargin={15mm}, hmargin={15mm, 15mm}}
\setlength{\parindent}{0pt}
\setlength{\parskip}{0.7\baselineskip}%{0.8em plus 1pt minus 1pt}

\onehalfspacing

% hyperref Settings
\hypersetup{
    colorlinks=true,
    linkcolor=blue,
    filecolor=magenta,      
    urlcolor=cyan,
}

% inkscape Settings
\newcommand{\incfig}[1]{%
    \def\svgwidth{\columnwidth}
    \import{./figure/}{#1.pdf_tex}
}

% inline code block settings
\newtcbox{\mybox}[1][]{
    on line,
    arc=1pt, outer arc=2pt,
    colback=lightgray!50!white, colframe=lightgray!50!white,
    boxsep=0pt, left=0pt, right=-1pt, top=2pt, bottom=0pt,
    boxrule=0pt, #1
}

\NewDocumentCommand{\code}{v}{%
  \begingroup\ttfamily
  \codeulinner{#1}%
  \endgroup%
}

% Definition Box Styles
\newtcbtheorem[auto counter, number within=subsection]{defn}{Definition}{
% basic settings
    theorem name and number,
    separator sign={},
    description delimiters={(}{)},
    enhanced jigsaw,
    sharp corners,
% box margins and rules
    boxrule=1pt,
    left=0.2cm,right=0.2cm,top=0.2cm,
    toptitle=0.1cm,
    bottomtitle=0.08cm,
% box colors
    colframe=white!25!black,
    colback=white,
    coltitle=white,
% box fonts
    fonttitle=\bfseries,fontupper=\normalsize
}{defn}

% Theorem Box Styles
\newtcbtheorem[auto counter, number within=subsection]{thm}{Theorem}{
% basic settings
    theorem name and number,
    separator sign={},
    description delimiters={(}{)},
    enhanced jigsaw,
    sharp corners,
% box margins and rules
    boxrule=1pt,
    left=0.2cm,right=0.2cm,top=0.2cm,
    toptitle=0.1cm,
    bottomtitle=0.08cm,
% box colors
    colframe=white!25!black,
    colback=white,
    coltitle=white,
% box fonts
    fonttitle=\bfseries,fontupper=\normalsize
}{thm}

% Lemma Box Styles
\newtcbtheorem[number within=tcb@cnt@thm]{lem}{Lemma}{
% basic settings
    theorem name and number,
    separator sign={},
    description delimiters={(}{)},
    enhanced jigsaw,
    sharp corners,
% box margins and rules
    boxrule=1pt,
    left=0.2cm,right=0.2cm,top=0.2cm,
    toptitle=0.1cm,
    bottomtitle=0.08cm,
% box colors
    colframe=white!50!black,
    colback=white,
    coltitle=white,
% box fonts
    fonttitle=\bfseries,fontupper=\normalsize
}{lem}

% Corollary Box Styles
\newtcbtheorem[number within=tcb@cnt@thm]{cor}{Corollary}{
% basic settings    
    theorem name and number,
    separator sign={},
    description delimiters={(}{)},
    enhanced jigsaw,
    sharp corners,
% box margins and rules
    boxrule=1pt,
    left=0.2cm,right=0.2cm,top=0.2cm,
    toptitle=0.1cm,
    bottomtitle=0.08cm,
% box colors
    colframe=white!50!black,
    colback=white,
    coltitle=white,
% box fonts
    fonttitle=\bfseries,fontupper=\normalsize
}{cor}

% Remark Box Styles
\newtcbtheorem[number within=subsection]{rem}{Remark}{
% basic settings
    theorem name and number,
    separator sign={},
    description delimiters={(}{)},
    enhanced jigsaw,
    sharp corners,
% box colors
    colframe=white!60!black,
    colback=white,
    coltitle=black,
% box margins and rules
    boxrule=1pt,
    left=0.2cm,right=0.2cm,top=0.2cm,
    toptitle=0.1cm,
    bottomtitle=0.08cm,
% box fonts
    fonttitle=\bfseries,fontupper=\normalsize
}{rem}

% Example Box Styles
% Remark Box Styles
\newtcbtheorem[number within=subsection]{exm}{Example}{
% basic settings
    theorem name and number,
    separator sign={},
    description delimiters={(}{)},
    enhanced jigsaw,
    sharp corners,
% box colors
    colframe=white!60!black,
    colback=white,
    coltitle=black,
% box margins and rules
    boxrule=1pt,
    left=0.2cm,right=0.2cm,top=0.2cm,
    toptitle=0.1cm,
    bottomtitle=0.08cm,
% box fonts
    fonttitle=\bfseries,fontupper=\normalsize
}{exm}


% Proof Box Styles

\title{Generative Adversarial Networks}
\author{Hyoung Joon, Kim}
\date{\today}

\begin{document}
\maketitle
\tableofcontents
\newpage

\section{Introduction}
A \textbf{generative adversarial network} (GAN) is an unsupervised network, which aims to
generate new samples that are indistinguishable from original data. GAN is just an architecture
for creating new samples, so it cannot evaluate probability if a new sample belongs to the original
distribution.


\section{Main Idea}
    \subsection{Generator and Discriminator}
        GAN uses two distinct networks; the \textbf{discriminator} and \textbf{generator}.
        The discriminator network classifies a given data; it classifies if the given data is created by generator or not.
        The generator network's objective is to create samples that can deceive a discriminator.
        The idea is simple, but GAN itself is hard to train. We'll discuss more about this in section \ref{sec:exGAN}.

        The generator creates new samples by first picking a latent variable $\textbf{z}$ from some base distribution
        and passing this variable through a network $\textbf{g}(\textbf{z}, \bm{\theta})$ with parameter $\bm{\theta}$.
        The goal is to find parameter $\bm{\theta}$ that creates sample indistinguishable from original dataset. This
        indistinguishableness, is defined to be statistically indistinguishable from the true data.

        The discriminator network $\textbf{f}(\textbf{x}, \bm{\phi})$ with parameter $\bm{\phi}$ discriminates
        whether the given data \textbf{x} is created by generator or not.

    \subsection{Training GAN}
        \subsubsection{Loss Function for GAN}
            The GAN uses binary cross entropy loss for training discriminator. Generator is trained to increase the loss, when
            discriminator parameter is fixed.
            \begin{defn}{Loss Function for GAN}{lossFnGAN}
                Let $\textbf{g}(\textbf{z}, \bm{\theta})$ and $\textbf{f}(\textbf{x}, \bm{\phi})$ be generator and discriminator, respectively,
                where $\textbf{z}_j$ is a latent variable and $\bm{\theta}$ and $\bm{\phi}$ are parameters.
                Let $\textbf{y}_j$ be a sample created by the generator with latent variable $\textbf{z}_j$,
                and let $\textbf{x}_i$ be the original data.
                Then, the loss for GAN is defined as following:
                \[
                \sum_{j} -\log{(1 - \sigma(\textbf{f}(\textbf{y}_j, \bm{\phi})))}
                -\sum_{i} -\log{(\theta(\textbf{f}(\textbf{x}_i, \bm{\phi})))}
                \]
                which can be written as
                \[
                \sum_{j} -\log{(1 - \sigma(\textbf{f}(\textbf{g}(\mathbf{z}_j, \bm{\theta}), \bm{\phi})))}
                -\sum_{i} -\log{(\theta(\textbf{f}(\textbf{x}_i, \bm{\phi})))}
                \]
                
            \end{defn}
        \subsubsection{Training Generator and Discriminator}
            Given the loss function above, the discriminator first tries to minimize the loss function,
            by adjusting $\bm{\phi}$. Then, the generator tries to maximize the loss with $\bm{\phi}$ fixed, by
            adjusting $\bm{\theta}$. Technically, the solution is \emph{Nash Equilibrium}.

            \begin{defn}{Training Discriminator}{trainDisc}
                The loss function for the discriminator is
                \[
                \mathcal L_{discriminator} (\bm{\phi}) = \sum_{j} -\log{(1 - \sigma(\textbf{f}(\textbf{g}(\mathbf{z}_j, \bm{\theta}), \bm{\phi})))}
                + \sum_{i} -\log{(\sigma(\textbf{f}(\textbf{x}_i, \bm{\phi})))}
                \]
                using the notation of Definition \ref{defn:lossFnGAN}.
            \end{defn}

            \begin{defn}{Training Generator}{trainGen}
                The loss function for the generator is
                \[
                \mathcal L_{generator} (\bm{\theta}) =
                \sum_{j} \log{(1 - \sigma(\textbf{f}(\textbf{g}(\mathbf{z}_j, \bm{\theta}), \bm{\phi})))}
                \]
                using the notation of Definition \ref{defn:lossFnGAN}. The second function is dropped since it has
                no dependence on $\bm{\theta}$, and the leftover is multiplied by -1.
            \end{defn}

            We draw a batch of latent variables $\textbf{z}_j$, and create samples. Then we choose a batch of real examples
            $\textbf{x}_i$. Then we perform gradient descents on each loss function.


\section{Extensions of GAN}\label{sec:exGAN}
    \subsection{WGAN}
        \subsubsection{Introduction}
            The original GAN architecture are difficult to train. The training is unstable itself,
            and even though it is trained well, it may suffer from \textbf{mode dropping} and \textbf{mode collapse}.
            If the generator succesfully deceives discriminator but creates only a subset of whole data, this is called mode dropping.
            If this phenomenon gets extreme, the generator creates only one sample, ignoring the latent variable.
            This is called mode collapse.

        \subsubsection{Analysis of GAN Loss Function}
            To figure out why the mode dropping and collapse happen, let's examine the loss function of
            GAN more precisely.
            \begin{thm}{Analysis of GAN Loss Function}{ganLossAnalysis}
                Let $p$ be the original distribution of the data, and let $q$ be the the distribution over the generated samples.
                Then, the loss of GAN function can be interpreted as the \textbf{Jensen-Shannon divergence}, ignoring the additive and multiplicative
                constants;
                \[
                \mathcal{L}_{discriminator}(\bm{\phi})=-\int q(\textbf{x})\log{\frac{q(\textbf{x})}{q(\textbf{x}) + p(\textbf{x})}}d\textbf{x}
                -\int p(\textbf{x})\log{\frac{p(\textbf{x})}{q(\textbf{x}) + p(\textbf{x})}}d\textbf{x}
                \]
                where Jensen-Shannon divergence is defined as:
                \[
                D_{JS}(q||p) = \frac{1}{2}\int q(x)\log{\frac{2q(\textbf{x})}{q(\textbf{x}) + p(\textbf{x})}}d\textbf{x}
                + \frac{1}{2}\int p(x)\log{\frac{2p(\textbf{x})}{q(\textbf{x}) + p(\textbf{x})}}d\textbf{x}
                \]
            \end{thm}

            Whenever the sample density $q(x)$ is high and the mixture $q(x) + p(x)$ has high probability, the first
            term of JS divergence will be small. That is, the first term penalizes regions where there generated samples are plenty
            but real samples are not; if enforces \emph{quality}.
            Vice versa, the second term penalizese regions with plenty of real samples and sparse generated samples; it enforces \emph{coverage}.
            Since the second term does not depend on generator, generator does not care about the coverage of the samples being created.
            This can be a reason of mode dropping and collapse.

        \subsubsection{Vanishing Gradients}
            When the discriminator is optimal, it maximizes the JS-divergence between two distributions.
            If the probability distributions are completely disjoint, this distance is maximal, which will be a constant of $\log{2}$.
            So any small change to the generator won't decrease the loss.
            The problem is that the distributions of samples and original dataset may really be disjoint, which
            makes training impossible.

        \subsubsection{EM Distance}
            To solve the problem, we use different distance metric called \textbf{EM distance} or \textbf{Wasserstein distance}.
            EM distance is defined as follows:
            \begin{defn}{EM Distance}{emDist}
                Let $p$ and $q$ be probability distribution. The \textbf{EM distance} $W(p,q)$ between $p$ and $q$ is defined by the equation
                \[
                W(p,q)=\inf_{\gamma \in \prod(p,q)}\mathbb E_{(x,y) \sim \gamma}[\|x-y\|]
                \]
                where $\prod(p,q)$ is sets of all possible joint distribution over $p$ and $q$.
            \end{defn}
            We can think of $\gamma$ as a transport plan. The EM distance, as we can see from its name,
            can be understood as the minmum transport cost used to transport $p$ to $q$.
            This can be understood better when the distributions are discrete.

            Using this equation directly is impossible since finding the infimum in
            high-dimensional space is tremendously costly. So, we use the duality property of the optimization theory.
            Then, the equation can be written as
            \[
            W(p,q)=\sup_{\|f\|_L \leq 1}(\mathbb E_{x \sim p}[f(x)] - \mathbb E_{y \sim q}[f(y)]),
            \]
            where $\|f\|_L \leq 1$ denotes the Lipschitz condition, $\|\nabla_{\textbf{x}}f(\textbf{x}, \bm{\phi})\| \leq 1$.
            Now the objective is to find the function $f$, called \textbf{critic}, which is usually trained with neural network.
        
        \subsubsection{WGAN}
            We now define Wasserstein GAN loss function and use it to train GAN.
            \begin{defn}{Wasserstein GAN Loss Function}{wassGANLossFn}
                Let $\textbf{z}_j$ be a latent variable and let $\textbf{x}_i$ be a given real samples.
                The \textbf{Wasserstein GAN loss function} is given as following:
                \[
                \mathcal L = \sum_{j} f(\textbf{g}(\textbf{z}_j, \bm{\theta}), \bm{\phi}) - \sum_{i} f(\textbf{x}_i, \bm{\phi}),
                \ \text{with} \ \|\nabla_{\textbf{x}}f(\textbf{x}, \bm{\phi})\| \leq 1
                \]
                where \textbf{g} is a generator network, and $f$ is a critic network.
            \end{defn}

            The expectation terms are approximated with sums about samples, and the condition $\|\nabla_{\textbf{x}}f(\textbf{x}, \bm{\phi})\| \leq 1$
            is called \textbf{Lipschitz condition}.

            One way the achieve this is to:
            \begin{enumerate}
                \item Clip the gradient; This is the what the original WGAN uses.
                \item Add a regularization term, penalizing when the gradient norm is bigger than one; This is called WGAN-GP,
                    where GP stands for \emph{gradient penalty}.
            \end{enumerate}

            The \emph{critic} has same objective as the discriminator of original GAN. However, the critic
            does not produce probability. Instead, it calculates a score that measures the sample's proximity to
            the real data distribution.

\end{document}